% !TEX root = ../main.tex

\section{第 15.1 节：两个自旋的合成与基的选择}
\label{section:QM-C15-S01}
\studymeta
  {QM-C15-S01}
  {2026-11-01 至 2026-11-02}
  {Shankar, \textit{Principles of Quantum Mechanics}, Second Edition}
  {\S 15.1 A Simple Example；印刷 pp.~403--408 / PDF pp.~409--414，至 \S 15.2 前；习题 15.1.2 只选第 (1) 问}
  {四块讲解与选定习题已核对；笔记已确认}

\notetopic{QM-C15-S01-T01}{总自旋与乘积基}
\keyidea{两个粒子的自旋可分别描述，也可按系统的总角动量分类。角动量合成要找出总角动量大小与一个分量的共同本征态。}

先忽略轨道自由度，考虑两个自旋 $1/2$ 的粒子。联合自旋空间及一组 product basis（乘积基）为
\begin{equation}
  \mathcal H_s=\complex^2\otimes\complex^2,\qquad
  \mathcal B_p=(\ket{++},\ket{+-},\ket{-+},\ket{--}).
\end{equation}
例如 $\ket{+-}=\ket+_1\otimes\ket-_2$，两个符号分别给出 $m_1,m_2=\pm1/2$。算符在联合空间上的作用为
\begin{equation}
  S_{1i}=\frac\hbar2\sigma_i\otimes I,\qquad
  S_{2i}=I\otimes\frac\hbar2\sigma_i,\qquad
  [S_{1i},S_{2j}]=0.
\end{equation}
因此乘积基同时对角化 $S_1^2,S_2^2,S_{1z},S_{2z}$，其中 $S_1^2=S_2^2=3\hbar^2I/4$。

定义总自旋
\begin{equation}
  \boxed{\mathbf S=\mathbf S_1+\mathbf S_2,\qquad S_z=S_{1z}+S_{2z}.}
\end{equation}
两个因子上的旋转相乘，给出共同旋转
\begin{equation}
  U(\mathbf n,\theta)
  =\ee^{-\ii\theta\mathbf n\cdot\mathbf S_1/\hbar}
   \ee^{-\ii\theta\mathbf n\cdot\mathbf S_2/\hbar}
  =\ee^{-\ii\theta\mathbf n\cdot\mathbf S/\hbar}.
\end{equation}
由各自的角动量代数及交叉对易性，也有 $[S_i,S_j]=\ii\hbar\sum_k\epsilon_{ijk}S_k$。所以 $\mathbf S$ 确实是整体旋转的生成元。

乘积基中的总投影为 $\hbar(m_1+m_2)$，即
\begin{equation}
  S_z\ \longleftrightarrow\ \hbar\diag(1,0,0,-1).
\end{equation}
$S_z=0$ 的本征空间由 $\ket{+-},\ket{-+}$ 张成；其中任意归一化叠加仍有确定的总投影，即使各自的 $S_{iz}$ 不确定。总自旋大小则由
\begin{equation}
  \boxed{S^2=S_1^2+S_2^2+2\mathbf S_1\cdot\mathbf S_2}
\end{equation}
决定。$S^2$ 与 $S_z$ 对易，但一般不与 $S_{1z},S_{2z}$ 对易，因而需要另一组共同本征基。

\notetopic{QM-C15-S01-T02}{总自旋基：三重态与单重态}
记 total-spin basis（总自旋基）为 $\ket{s,m}$，省略固定的 $s_1=s_2=1/2$。其含义是
\begin{equation}
  S^2\ket{s,m}=s(s+1)\hbar^2\ket{s,m},\qquad
  S_z\ket{s,m}=m\hbar\ket{s,m}.
\end{equation}
$m=1$ 子空间只有 $\ket{++}$ 一个基向量。由于 $[S^2,S_z]=0$，它也是 $S^2$ 的本征态；再由 $S_+\ket{++}=0$，确定它是 $s=1$ 的最高态。总降算符为 $S_-=S_{1-}+S_{2-}$，故
\begin{equation}
  \hbar\sqrt2\ket{1,0}=S_-\ket{1,1}
    =\hbar(\ket{-+}+\ket{+-}).
\end{equation}
再次下降得到 $\ket{1,-1}=\ket{--}$。剩余正交组合是反对称的 $m=0$ 态；它被 $S_z,S_+,S_-$ 全部消去，由 $S^2=S_+S_-+S_z^2-\hbar S_z$ 得到 $s=0$。

\begin{center}
\renewcommand{\arraystretch}{1.25}
\begin{tabular}{@{}llll@{}}
  \toprule
  总自旋基 & 乘积基表达 & $S^2$ & $S_z$ \\
  \midrule
  $\ket{1,1}$ & $\ket{++}$ & $2\hbar^2$ & $\hbar$ \\
  $\ket{1,0}$ & $(\ket{+-}+\ket{-+})/\sqrt2$ & $2\hbar^2$ & $0$ \\
  $\ket{1,-1}$ & $\ket{--}$ & $2\hbar^2$ & $-\hbar$ \\
  $\ket{0,0}$ & $(\ket{+-}-\ket{-+})/\sqrt2$ & $0$ & $0$ \\
  \bottomrule
\end{tabular}
\end{center}
三个 $s=1$ 态称为 triplet（三重态），一个 $s=0$ 态称为 singlet（单重态）。$s=1$ 的总角动量大小为 $\sqrt2\hbar$。这些名称先指态空间的维数；能量与能级简并还取决于 Hamiltonian。矩阵验证见\relatedexercise{QM-C15-S01-E02}{习题 15.1.1}。

两套基描述同一个四维空间，换基保留物理状态。空间分解简记为
\begin{equation}
  \boxed{\tfrac12\otimes\tfrac12=1\oplus0,\qquad
    2\times2=3+1.}
\end{equation}
左边表示两个自旋空间的张量积，右边表示总自旋 $s=1,0$ 子空间的直和。因 $[S^2,S_i]=0$，共同旋转保持这两个子空间；在总自旋基中，旋转矩阵分成 $3\times3$ 与 $1\times1$ 两块。所有方向的旋转共同作用时，三维的 $s=1$ 空间不能再拆成更小的不变子空间，这称为不可约；单重态满足 $\mathbf S\ket{0,0}=0$，在共同旋转下不变。

\begin{center}
\begin{tikzpicture}[x=1cm,y=1cm,>=Stealth,font=\small]
  \node[noteBlue] at (2.8,0.75) {$s=1$：三维子空间};
  \node (a) at (0,0) {$\ket{1,1}$};
  \node (b) at (2.8,0) {$\ket{1,0}$};
  \node (c) at (5.6,0) {$\ket{1,-1}$};
  \draw[->,thick,noteBlue] (a)--node[above] {$S_-$} node[below] {$\hbar\sqrt2$} (b);
  \draw[->,thick,noteBlue] (b)--node[above] {$S_-$} node[below] {$\hbar\sqrt2$} (c);
  \draw[dashed,gray] (7,-0.6)--(7,1.0);
  \node[accentRed] at (9.2,0.75) {$s=0$：一维子空间};
  \node at (9.2,0) {$\ket{0,0}$};
  \node at (9.2,-0.65) {$S_\pm\ket{0,0}=0$};
\end{tikzpicture}

{\small 原创结构示意：箭头表示降算符作用，不表示时间演化或能级高低。}
\end{center}

\notetopic{QM-C15-S01-T03}{交换对称性与全同费米子}
交换两个自旋标签使 $\ket{+-}\leftrightarrow\ket{-+}$，而 $\ket{++},\ket{--}$ 保持不变。因此三重态的自旋部分对称，单重态的自旋部分反对称。这里交换的是粒子标签；宇称则反演空间坐标。

对两个全同费米子，交换位置和自旋的全部标签时，\emph{完整态}必须反对称。若态可分离为 $\Psi=\Phi_{\rm space}\chi_{\rm spin}$，空间与自旋的交换符号须满足 $\eta_o\eta_s=-1$：
\begin{center}
\begin{tabular}{@{}lll@{}}
  \toprule
  空间部分 & 自旋部分 & 完整态的交换符号 \\
  \midrule
  对称 & 单重态（反对称） & $-1$ \\
  反对称 & 三重态（对称） & $-1$ \\
  \bottomrule
\end{tabular}
\end{center}
一般允许态可以由这些乘积态线性组合而成。对两个正交归一的不同空间轨道 $u,v$，
\begin{equation}
  \Phi_\pm(\mathbf r_1,\mathbf r_2)
  =\frac{u(\mathbf r_1)v(\mathbf r_2)\pm v(\mathbf r_1)u(\mathbf r_2)}{\sqrt2}
\end{equation}
分别与单重态、三重态配合。若两电子占据同一空间轨道，空间乘积已对称，反对称的空间组合则为零，因此只能配单重态。氦原子 $1s^2$ 的教材示范见\relatedexercise{QM-C15-S01-E01}{完整态的构造}。

\notetopic{QM-C15-S01-T04}{Hamiltonian 决定方便的基}
\keyidea{优先选使 Hamiltonian 对角化的基：分别耦合各自投影时用乘积基，耦合总自旋大小时用总自旋基。}

\textbf{均匀外磁场。}两个自旋分别与 $\mathbf B=B_0\hat{\mathbf z}$ 耦合、无直接自旋间耦合时，
\begin{equation}
  H=-B_0(\gamma_1S_{1z}+\gamma_2S_{2z}),\qquad
  \boxed{E_{m_1,m_2}=-B_0\hbar(\gamma_1m_1+\gamma_2m_2).}
\end{equation}
乘积基是能量本征基。若 $\gamma_1=\gamma_2=\gamma$，则 $H=-\gamma B_0S_z$，总自旋基也能对角化它。

\textbf{各向同性自旋耦合。}对实常数 $A$，
\begin{equation}
  \boxed{H=A\mathbf S_1\cdot\mathbf S_2
     =\frac A2(S^2-S_1^2-S_2^2)
     =\frac A2S^2-\frac34A\hbar^2I.}
\end{equation}
此时总自旋基对角化 $H$，能量仅依赖 $s$；三个三重态同能。$A$ 的量纲为能量除以角动量平方。

\textbf{超精细相互作用。}电子与质子的磁矩相互作用称为 hyperfine interaction（超精细相互作用）。教材将氢原子空间轨道固定为 $1s$，采用有效模型
\begin{equation}
  H=-\mathrm{Ry}\,I+A\mathbf S_e\cdot\mathbf S_p,\qquad A>0.
\end{equation}
$\mathrm{Ry}$ 为里德伯能量。电子、质子各有自旋 $1/2$，且是可区分粒子，单重态和三重态均可存在。耦合将同一库仑能量分裂为两档，计算见\relatedexercise{QM-C15-S01-E03}{习题 15.1.2 第 (1) 问}。

\textbf{一分钟回顾。}乘积基确定各自的 $z$ 分量，总自旋基确定整体的 $S^2,S_z$。两个自旋 $1/2$ 合成一个三重态子空间和一个单重态子空间，共同旋转保持二者。全同费米子的空间与自旋部分须配成完整反对称态。选择基时看 Hamiltonian：$S_{1z},S_{2z}$ 型作用适合乘积基，$\mathbf S_1\cdot\mathbf S_2$ 型作用适合总自旋基。

\section{第 15.2 节：一般角动量合成与 CG 系数}
\label{section:QM-C15-S02}
\studymeta
  {QM-C15-S02}
  {2026-11-02 至 2026-11-04}
  {Shankar, \textit{Principles of Quantum Mechanics}, Second Edition}
  {\S 15.2 The General Problem；印刷 pp.~408--416 / PDF pp.~414--422，至 \S 15.3 标题前}
  {四块讲解与两道选定习题已核对；笔记已确认}

\notetopic{QM-C15-S02-T01}{一般合成规则与两套基}
\keyidea{固定两个角动量量子数 $j_1,j_2$，把各自分量确定的乘积基，换成总角动量 $J^2,J_z$ 确定的耦合基。}

设 $\mathbf J_1,\mathbf J_2$ 各自满足角动量代数，且 $[J_{1i},J_{2k}]=0$。它们可对应两个粒子，也可对应同一粒子的轨道与自旋。总角动量为
\begin{equation}
  \mathbf J=\mathbf J_1+\mathbf J_2,\qquad
  J_z=J_{1z}+J_{2z},\qquad J_\pm=J_{1\pm}+J_{2\pm}.
\end{equation}
以下固定 $J_i^2=j_i(j_i+1)\hbar^2I$，在 $(2j_1+1)(2j_2+1)$ 维空间内讨论。
\begin{center}
\begin{tabularx}{\textwidth}{@{}lXX@{}}
  \toprule
  基 & 基向量 & 另外确定的观测量 \\
  \midrule
  乘积基 & $\ket{j_1,m_1}\otimes\ket{j_2,m_2}$ & $J_{1z},J_{2z}$ \\
  耦合基 & $\ket{j,m;j_1,j_2}\equiv\ket{j,m}$ & $J^2,J_z$ \\
  \bottomrule
\end{tabularx}
\end{center}
这里省略 $\ket{j,m}$ 中固定的 $j_1,j_2$，并有
\begin{equation}
  J^2\ket{j,m}=j(j+1)\hbar^2\ket{j,m},\qquad
  J_z\ket{j,m}=m\hbar\ket{j,m}.
\end{equation}
乘积态的总投影是 $m=m_1+m_2$；同一 $m$ 往往对应多个乘积态，须在这个子空间内进一步对角化 $J^2$。

一般合成规则为
\begin{equation}
  \boxed{j=|j_1-j_2|,\ |j_1-j_2|+1,\ldots,j_1+j_2,\qquad
    m=-j,-j+1,\ldots,j.}
\end{equation}
$j$ 的步长为 $1$，且与 $j_1+j_2$ 同为整数或半整数。每个允许的 $j$ 在这个固定张量积中出现一次：
\begin{equation}
  \boxed{\mathcal H_{j_1}\otimes\mathcal H_{j_2}
    =\bigoplus_{j=|j_1-j_2|}^{j_1+j_2}\mathcal H_j.}
\end{equation}
每个 $\mathcal H_j$ 含 $2j+1$ 个态，在共同旋转下保持不变。令 $j_{\max}=j_1+j_2$、$j_{\min}=|j_1-j_2|$，维数核对为
\begin{equation}
  \sum_j(2j+1)
  =(j_{\max}-j_{\min}+1)(j_{\max}+j_{\min}+1)
  =(2j_1+1)(2j_2+1).
\end{equation}
教材印刷 p.~409 先用维数检验提出的分解。严格地说，数目相同只检验计数；各个子空间的存在由下面的最高态构造确定。

\notetopic{QM-C15-S02-T02}{最高态构造与 Clebsch--Gordan 系数}
\textbf{先构造最大 $j$。}最大投影 $m=j_{\max}$ 只有一个乘积态，因此
\begin{equation}
  \boxed{\ket{j_{\max},j_{\max}}
    =\ket{j_1,j_1}\otimes\ket{j_2,j_2}.}
\end{equation}
它被 $J_+$ 消去，由 $J^2=J_-J_++J_z^2+\hbar J_z$ 确定 $j=j_{\max}$。用
\begin{equation}
  J_-\ket{j,m}=\hbar\sqrt{(j+m)(j-m+1)}\ket{j,m-1}
\end{equation}
逐次下降，即可生成该 $j$ 的全部态。若 $j_1,j_2>0$，记
\begin{equation}
  \ket a=\ket{j_1,j_1-1}\otimes\ket{j_2,j_2},\qquad
  \ket b=\ket{j_1,j_1}\otimes\ket{j_2,j_2-1}.
\end{equation}
比较 $J_-=J_{1-}+J_{2-}$ 在两套基中的作用，得到
\begin{equation}
  \boxed{\ket{j_{\max},j_{\max}-1}
    =\sqrt{\frac{j_1}{j_{\max}}}\ket a
     +\sqrt{\frac{j_2}{j_{\max}}}\ket b.}
\end{equation}

\textbf{再找新的最高态。}在同一 $m=j_{\max}-1$ 子空间中，取归一化的正交组合
\begin{equation}
  \ket\chi=\sqrt{\frac{j_1}{j_{\max}}}\ket b
           -\sqrt{\frac{j_2}{j_{\max}}}\ket a.
\end{equation}
$J_+$ 作用于 $\ket b,\ket a$ 的系数分别为 $\hbar\sqrt{2j_2},\hbar\sqrt{2j_1}$，两项恰好抵消。故 $J_+\ket\chi=0$，结合 $m=j_{\max}-1$，有
\begin{equation}
  \boxed{\ket\chi=\ket{j_{\max}-1,j_{\max}-1}.}
\end{equation}
对它继续下降，再在尚未覆盖的正交补中寻找最高态，直至 $j_{\min}$。这一过程逐个构造各 $j$ 子空间；若某个 $j_i=0$，则只有另一个 $j$，无需上述两项组合。

\textbf{升降只改变 $m$。}由于 $[J^2,J_\pm]=0$，$J_+$ 与 $J_-$ 均保留 $j$。生成另一组 $j$ 的态，必须在正交补中另选最高态。三重态与单重态的投影形式见\relatedexercise{QM-C15-S02-E03}{习题 15.2.5}。

\textbf{CG 系数就是换基振幅。}定义 Clebsch--Gordan coefficients（Clebsch--Gordan 系数，简称 CG 系数）
\begin{equation}
  C^{jm}_{j_1m_1;j_2m_2}
    =\bigl(\bra{j_1,m_1}\otimes\bra{j_2,m_2}\bigr)\ket{j,m},
\end{equation}
从而
\begin{equation}
  \boxed{\ket{j,m}
    =\sum_{m_1,m_2}C^{jm}_{j_1m_1;j_2m_2}
      \ket{j_1,m_1}\otimes\ket{j_2,m_2}.}
\end{equation}
系数只有在 $j$ 满足合成规则且 $m_1+m_2=m$ 时才可能非零。在 $\ket{j,m}$ 中联合测量 $J_{1z},J_{2z}$，结果 $(m_1\hbar,m_2\hbar)$ 的概率就是 $|C^{jm}_{j_1m_1;j_2m_2}|^2$。

教材的相位约定取所有系数为实数，并使每组最高态 $\ket{j,j}$ 中 $m_1=j_1$ 的乘积项系数为正；同一组其余态的相位由正的下降系数固定。在这一约定下，
\begin{equation}
  \boxed{C^{j,-m}_{j_1,-m_1;j_2,-m_2}
    =(-1)^{j_1+j_2-j}C^{jm}_{j_1m_1;j_2m_2}.}
\end{equation}
这可直接补齐负 $m$ 的系数。两套基均正交归一，因此 CG 矩阵是酉矩阵；实系数约定下为正交矩阵。逆变换为
\begin{equation}
  \ket{j_1,m_1}\otimes\ket{j_2,m_2}
    =\sum_{j,m}\bigl(C^{jm}_{j_1m_1;j_2m_2}\bigr)^*\ket{j,m}.
\end{equation}
矩阵示范见\relatedexercise{QM-C15-S02-E01}{教材正文换基示范}；完整的 $\tfrac12\otimes1$ 构造见\relatedexercise{QM-C15-S02-E02}{习题 15.2.2 第 (1) 问}。

\notetopic{QM-C15-S02-T03}{轨道角动量与自旋的合成}
对电子 $s=1/2$，固定轨道量子数 $l$，有
\begin{equation}
  \boxed{\mathbf J=\mathbf L+\mathbf S,\qquad
    j=l+\tfrac12\ \text{或}\ l-\tfrac12.}
\end{equation}
下支只在 $l\ge1$ 时存在；$l=0$ 时只有 $j=1/2$。这里的两支由角动量代数决定，不是经典矢量严格平行或反平行的假设。

固定总投影 $m$，在两项均存在的区间 $|m|\le l-1/2$，定义
\begin{equation}
  \ket u=\ket{l,m-\tfrac12}\otimes\ket+,
  \qquad
  \ket v=\ket{l,m+\tfrac12}\otimes\ket-.
\end{equation}
由 $J^2=L^2+S^2+2L_zS_z+L_-S_++L_+S_-$，在 $(\ket u,\ket v)$ 中
\begin{equation}
  \frac{J^2}{\hbar^2}\ \longleftrightarrow\
  \begin{pmatrix}
    l(l+1)+m+\tfrac14&q\\
    q&l(l+1)-m+\tfrac14
  \end{pmatrix},\qquad
  q=\sqrt{(l+\tfrac12)^2-m^2}.
\end{equation}
令上支本征态为 $\alpha\ket u+\beta\ket v$，比较第一行得到
\begin{equation}
  (l+\tfrac12-m)\alpha=q\beta,
  \qquad \frac\beta\alpha
    =\sqrt{\frac{l+\tfrac12-m}{l+\tfrac12+m}}.
\end{equation}
结合归一化、正交性及教材相位约定，记
\begin{equation}
  A_m=\sqrt{\frac{l+\tfrac12+m}{2l+1}},\qquad
  B_m=\sqrt{\frac{l+\tfrac12-m}{2l+1}},
\end{equation}
得到
\begin{equation}
  \boxed{\begin{aligned}
    \ket{l+\tfrac12,m}&=A_m\ket u+B_m\ket v,\\
    \ket{l-\tfrac12,m}&=-B_m\ket u+A_m\ket v.
  \end{aligned}}
\end{equation}
在上支的两个端点，只有一个合法乘积态：
\begin{equation}
  \ket{l+\tfrac12,l+\tfrac12}=\ket{l,l}\otimes\ket+,
  \qquad
  \ket{l+\tfrac12,-l-\tfrac12}=\ket{l,-l}\otimes\ket-.
\end{equation}
端点以外，两项通常均非零：$J^2,J_z$ 确定，而各自的 $L_z,S_z$ 具有非零涨落。

耦合基尤其适合 spin--orbit interaction（自旋--轨道相互作用），因为
\begin{equation}
  \mathbf L\cdot\mathbf S=\frac12(J^2-L^2-S^2),\qquad
  \boxed{\mathbf L\cdot\mathbf S\ \longleftrightarrow\
  \begin{cases}
    \hbar^2l/2,&j=l+1/2,\\
    -\hbar^2(l+1)/2,&j=l-1/2.
  \end{cases}}
\end{equation}
这说明含 $\mathbf L\cdot\mathbf S$ 的相互作用为什么适合用总角动量分类。具体能量分裂还需给出耦合强度及径向部分。

\notetopic{QM-C15-S02-T04}{光谱项符号与磁简并}
光谱中用 term symbol（光谱项符号）
\begin{equation}
  \boxed{{}^{2S+1}L_J}
\end{equation}
标记总自旋、总轨道角动量和合成后的总角动量。此处 $L,S,J$ 是量子数；多电子原子中指所有电子合成后的总量子数。
\begin{center}
\begin{tabularx}{\textwidth}{@{}lX@{}}
  \toprule
  位置 & 含义 \\
  \midrule
  左上角 $2S+1$ & 自旋多重度，如 $S=0$ 为单重、$S=1$ 为三重 \\
  中间字母 $L$ & 用 $\mathrm S,\mathrm P,\mathrm D,\mathrm F,\mathrm G,\ldots$ 表示 $L=0,1,2,3,4,\ldots$ \\
  右下角 $J$ & 总角动量量子数，$\mathbf J=\mathbf L+\mathbf S$ \\
  \bottomrule
\end{tabularx}
\end{center}
中间的字母 $\mathrm S$ 表示 $L=0$，与左上角公式中的自旋量子数 $S$ 是不同含义。例如
\begin{equation}
  {}^2\mathrm P_{3/2}:\ L=1,\ S=\tfrac12,\ J=\tfrac32,
  \qquad
  {}^2\mathrm P_{1/2}:\ L=1,\ S=\tfrac12,\ J=\tfrac12.
\end{equation}
氦原子基态的记号 $ {}^1\mathrm S_0$ 对应 $L=S=J=0$。

项符号没有给出磁量子数 $M_J$ 或完整的径向、组态信息。固定 $J$ 的多重态有 $M_J=-J,\ldots,J$ 共 $2J+1$ 个态；在旋转不变的 Hamiltonian 下，这些磁子态同能，外场可使其分裂。$2S+1$ 与 $2J+1$ 分别计数自旋多重度和总角动量磁子态。

\textbf{一分钟回顾。}两个角动量合成时，$j$ 从 $|j_1-j_2|$ 到 $j_1+j_2$ 按 $1$ 递增。先找最高态，再下降生成固定 $j$ 的一组态；新的 $j$ 要从正交补中构造。CG 系数是两套正交基间的变换振幅，模平方给出分量测量概率。电子的 $L+S$ 分成 $j=l\pm1/2$ 两支，光谱项符号 $ {}^{2S+1}L_J$ 汇总这些量子数，但不指定 $M_J$。
